\section{Research agenda and present status}\label{sec:outlook}

The revised program starts from a known arithmetic evolution and
asks for a physical derivation of its transfer and norm. The
detailed work remains useful because it tests specific candidates
for that derivation. The immediate priority is to define the
candidate boundary operator precisely enough that its variation
can be compared with \eqref{eq:physical-flow}.

\subsection{The first discriminating tasks}

First fix one field and endpoint construction: input/output
identifications, reference color sector, regulator, adjoint and
normalization, together with the meaning of $\omega$.
The existing free endpoint kernel is a control for nonzero pairing;
a mere family of scalar expectations is not yet the required
operator.

Next compute its effective insertion on a controlled sector.
The exact Wilson variation supplies the terms, and
\eqref{eq:generator} supplies the comparison. Record separately
whether the local constant, gamma difference, pole terms and
prime-power delays have been obtained. A missing component should
remain an explicit obstruction rather than be absorbed into an
unspecified normalization.

Then test the cumulative identity
\eqref{eq:norm-hypothesis}, or a physical form of the cross-interval
inequality \eqref{eq:gluing-main}. For the current bulge candidate,
the complete first interaction-order response is still relevant,
including endpoint, defect and junction contributions. Its purpose
is now to test closure and the norm balance. The detailed list of
missing diagrams and the nonuniform flat-limit issue are retained
in \cite[Section~S8.4]{Supplement}.

Finally formulate one rigorous joint depth--shift step and the
invariant needed to repeat it to unbounded depth. The new
\emph{critical-path} investigation is the place to develop that
continuation, starting from the inherited storage criteria
\cite{ShiftFlowNote,CumulativePath}. No such step or all-depth
bound is claimed by this restructuring.

\subsection{A compact statement ledger}

\renewcommand{\arraystretch}{1.13}
\begin{longtable}{@{}p{0.36\textwidth}p{0.59\textwidth}@{}}
\toprule
Ingredient & Present status\\
\midrule
\endfirsthead
\toprule
Ingredient & Present status (continued)\\
\midrule
\endhead
\bottomrule
\endfoot
Arithmetic shift and defect ODEs &
Exact inherited identities, with their causal and form domains.\\
Weil limit and diagonal implication &
Background positivity criterion and a sufficient conditional route
to RH; the required contractions remain to be proved.\\
Free tower and gamma model &
Explicit spectral ingredients with a fixed subtraction sign;
no complete transfer or physical mode selection.\\
Contour and Wilson variation tools &
Exact smooth identities and scoped rigidity/closure restrictions;
no arithmetic driver construction.\\
Chiral projectors and scalar endpoints &
Compatible charges exist; the ordinary same-charge scalar
pairing vanishes in the stated ansatz.\\
Ordinary reflected pairing &
A nonzero positive free kernel; interacting positivity and its
identification with cumulative storage remain open.\\
Smooth bulge response &
A nonzero direct Gaussian bulk sector; the full interacting
response and protection question remain open.\\
Physical transfer and positive balance &
Research hypotheses, including exact generator closure,
normalization, uniqueness and an independent amplitude map.\\
Depth--shift gluing &
Exact conditional Schur criteria; their all-input estimates
and an unbounded continuation remain open.\\
\end{longtable}
\renewcommand{\arraystretch}{1}

The companion preserves the technical evidence behind this ledger,
including the 268 finite diagnostics. Their scope is signs, ranks,
ordered variations and Gaussian integrals. They neither verify the
new realization hypothesis nor certify interacting positivity or RH.
